\subsection{Main Contribution}

The main contribution of this study is not the discovery of a complete swarm
mechanism. It is the isolation of a bounded local residual that survives
several natural reductions. This distinction matters because weakly ordered
midge swarms can contain reproducible organization even when global alignment
and simple low-dimensional closures are insufficient.

The positive and negative results work together. The positive evidence shows
that a local tangential non-affine residual is common after local affine
subtraction. The negative evidence shows that several simple explanatory
routes do not absorb this residual under their tested definitions. The
remaining object is therefore more precise: it is local, non-affine,
tangential, transition-linked, common in most observations, and explicitly
heterogeneous.

This interpretation is consistent with previous descriptions of midge swarms
as weakly ordered nonequilibrium collectives with spatial structure,
stochastic components, and long-range or effective interactions
\cite{feng2023_sampling,reynolds2017_velocity_gravity,reynolds2018_langevin,reynolds2021_intrinsic,reynolds2024_spatial,sinhuber2021_eos}.
The present analysis adds a reduction-based boundary: a reproducible local
residual remains after affine motion is removed, but it should not yet be
identified with a universal low-dimensional mechanism.

